%% ========================================================================
%% Bab 1: Arsitektur Web Modern dan Ekosistem Laravel
%% ========================================================================

\chapter{Arsitektur Web Modern dan Ekosistem Laravel}
\label{ch:arsitektur-web-modern}

\chapterhero{blue}{Bandingkan arsitektur monolith, microservices, dan serverless, lalu siapkan proyek Laravel 13 dengan SQLite sebagai fondasi Simple POS.}{\faIcon{server}\quad monolith vs microservices}{\faIcon{code-branch}\quad ekosistem Laravel}{\faIcon{database}\quad Simple POS}

Sebuah restoran keluarga yang baru buka biasanya menangani semuanya sendiri:
pemilik memasak, anaknya mencatat pesanan, istrinya menghitung uang di meja
kasir yang sama dengan meja pelanggan makan. Semua orang tahu semua hal, dan
kalau dapur kebanjiran pesanan pada malam Minggu, solusinya cukup menambah
kompor, bukan membuka cabang baru di ujung kota. Bandingkan dengan food
court di sebuah mal, tempat setiap tenant punya dapur sendiri, kasir
sendiri, bahkan resep rahasia sendiri, dan pengelola mal hanya menyediakan
meja bersama serta jalur listrik. Kalau tenant bakso ramai sementara tenant
es krim sepi, keduanya tetap berjalan sebagai unit terpisah, tidak saling
menunggu.

Kedua model itu punya padanan persis dalam arsitektur perangkat lunak:
monolith seperti restoran keluarga (satu basis kode, satu proses, satu tim
yang mengurus semuanya), dan microservices seperti food court (banyak
layanan independen yang saling terhubung lewat jaringan). Simple POS,
aplikasi kasir yang dibangun bab demi bab sepanjang buku ini, dipilih
sebagai monolith bukan karena keterbatasan, melainkan karena skalanya
cocok: sebuah warung dengan satu atau dua kasir tidak butuh sepuluh layanan
terpisah hanya untuk mencatat penjualan segelas kopi. Salah pilih arsitektur
di titik ini, membangun food court untuk bisnis yang cukup dilayani satu
dapur, berarti menghabiskan waktu pengembangan untuk mengurus jaringan
antar-layanan yang sebenarnya tidak dibutuhkan, alih-alih menyelesaikan
fitur yang dipakai kasir setiap hari.

\textbf{Yang Akan Kamu Pelajari}

\begin{enumerate}
  \item membandingkan arsitektur monolith, microservices, dan serverless untuk menjelaskan alasan Laravel dipilih sebagai kerangka kerja monolith bagi Simple POS;
  \item menyiapkan proyek Laravel 13 baru dengan SQLite sebagai basis data zero-setup, lengkap dengan \artisan{migrate:fresh --seed} dan \artisan{serve};
  \item menjelaskan konvensi commit \cmd{increment N} pada repositori Simple POS dan bagaimana urutan increment memetakan pertumbuhan aplikasi bab demi bab.
\end{enumerate}

\section{Monolith, Microservices, dan Mengapa Laravel}
\label{sec:arsitektur-web-modern-1}

\begin{istilahpenting}
\begin{description}[leftmargin=0pt,style=nextline]
  \item[Monolith] Aplikasi yang seluruh lapisannya, antarmuka, logika bisnis, dan akses data, berjalan dalam satu basis kode dan satu proses deploy.
  \item[Microservices] Arsitektur yang memecah aplikasi menjadi layanan-layanan independen, masing-masing berjalan dan di-deploy sendiri, saling berkomunikasi lewat jaringan.
  \item[MVC] Pola \textit{Model-View-Controller}: Model mengurus data, View mengurus tampilan, Controller mengurus alur permintaan di antara keduanya.
  \item[Artisan] Antarmuka baris perintah bawaan Laravel untuk menjalankan tugas pengembangan sehari-hari: migrasi, seeding, dan pembuatan berkas boilerplate.
  \item[Migrasi] Berkas PHP yang mendefinisikan perubahan skema basis data secara terprogram, sehingga skema dapat dibangun ulang secara konsisten di komputer mana pun.
\end{description}
\end{istilahpenting}

Sebuah aplikasi \term{monolith} menyimpan seluruh kodenya, mulai dari
tampilan yang dilihat pengguna sampai query yang menyentuh basis data, dalam
satu proyek yang di-deploy sebagai satu unit. Menambah fitur berarti
menambah kode pada proyek yang sama; men-deploy pembaruan berarti mengganti
satu unit itu dengan versi barunya. Sebaliknya, \term{microservices}
memecah tanggung jawab itu ke layanan-layanan terpisah, misalnya satu
layanan khusus mengurus katalog produk, layanan lain khusus mengurus
pembayaran, dan keduanya saling bicara lewat panggilan jaringan (biasanya
HTTP atau message queue). Pemisahan ini memberi kebebasan: setiap layanan
bisa ditulis dengan bahasa berbeda, di-deploy dengan jadwal berbeda, dan
diskalakan sendiri-sendiri saat trafiknya melonjak.

Kebebasan itu punya harga. Sebuah tim yang memecah aplikasinya menjadi
delapan layanan terpisah kini harus mengurus delapan proses deploy, delapan
titik yang bisa gagal, dan komunikasi jaringan antar-layanan yang sebelumnya
cukup dipanggil sebagai satu fungsi biasa di dalam proses yang sama. Untuk
tim besar yang mengoperasikan sistem berskala jutaan pengguna, ongkos itu
sepadan dengan kebebasan yang didapat. Untuk aplikasi seperti Simple POS,
yang melayani satu warung dengan volume transaksi harian yang bisa dihitung
manual di kertas, ongkos itu jauh melebihi manfaatnya.

\begin{table}[htbp]
\centering
\caption{Perbandingan arsitektur monolith, microservices, dan serverless}
\label{tab:arsitektur-perbandingan}
\begin{tblr}{
  colspec = {X[1,l] X[1.2,l] X[1,l] X[1.4,l]},
  hline{1,2,Z} = {solid},
}
Arsitektur & Deployment & Kompleksitas Awal & Cocok Untuk \\
Monolith & Satu unit & Rendah & Simple POS, MVP, tim kecil \\
Microservices & Banyak unit independen & Tinggi & Sistem skala besar, tim besar \\
Serverless & Fungsi per event & Sedang & Beban kerja sporadis \\
\end{tblr}
\end{table}

\begin{figure}[htbp]
\centering
\resizebox{0.92\linewidth}{!}{%
\begin{tikzpicture}[node distance=3mm]
  \node[diagbox accent, minimum width=3.2cm] (ui) {Antarmuka (UI)};
  \node[diagbox accent, minimum width=3.2cm, below=of ui] (logic) {Logika Bisnis};
  \node[diagbox accent, minimum width=3.2cm, below=of logic] (data) {Akses Data};
  \node[draw=black!45, thick, rounded corners=3pt, fit=(ui)(logic)(data), inner sep=7pt] {};
  \node[above=2mm of ui, font=\bfseries] {Monolith};

  \node[diagbox, right=3.4cm of ui, yshift=-0.3cm] (svc1) {Layanan Produk};
  \node[diagbox, below=1.4cm of svc1] (svc2) {Layanan Pembayaran};
  \node[diagbox, right=1.6cm of svc1, yshift=-0.7cm] (svc3) {Layanan Pengguna};
  \node[above=2mm of svc1, font=\bfseries] {Microservices};

  \draw[diagarrow, <->] (svc1) -- (svc2);
  \draw[diagarrow, <->] (svc2) -- (svc3);
  \draw[diagarrow, <->] (svc3) -- (svc1);
\end{tikzpicture}%
}
\caption{Struktur monolith (satu basis kode, tiga lapisan menyatu) vs
microservices (layanan independen yang saling memanggil lewat jaringan)}
\label{fig:monolith-microservices}
\end{figure}

\term{Serverless} menempuh jalan ketiga: kode ditulis sebagai fungsi-fungsi
kecil yang hanya berjalan ketika ada peristiwa yang memicunya (permintaan
HTTP, jadwal cron, pesan baru di antrean), dan penyedia layanan cloud yang
mengurus server di baliknya. Cocok untuk beban kerja yang naik-turun tajam,
misalnya pemrosesan gambar yang hanya dipicu saat pengguna mengunggah foto,
tetapi kurang cocok untuk aplikasi seperti Simple POS yang butuh koneksi
basis data yang konsisten sepanjang jam operasional warung.

Laravel dipilih sebagai kerangka kerja buku ini karena memetakan tepat ke
kebutuhan monolith: satu proyek PHP yang sudah menyediakan routing,
autentikasi, ORM, dan sistem template tanpa perlu merakit potongan pustaka
dari sumber berbeda. Dibandingkan menulis PHP polos dari nol, Laravel
memberi struktur MVC yang konsisten sejak baris kode pertama; dibandingkan
kerangka kerja microservices-first yang lebih rumit, Laravel tetap
produktif untuk tim satu atau dua orang yang mengejar tenggat.

\begin{figure}[htbp]
\centering
\resizebox{\linewidth}{!}{%
\begin{tikzpicture}[node distance=5mm]
  \node[diagbox accent] (req) {Request};
  \node[diagbox accent, right=of req] (router) {Router};
  \node[diagbox accent, right=of router] (ctrl) {Controller};
  \node[diagbox accent, right=of ctrl] (model) {Model};
  \node[diagbox accent, right=of model] (view) {View /\\Response};
  \draw[diagarrow] (req) -- (router);
  \draw[diagarrow] (router) -- (ctrl);
  \draw[diagarrow] (ctrl) -- (model);
  \draw[diagarrow] (model) -- (view);
\end{tikzpicture}%
}
\caption{Alur satu request Laravel: dari permintaan HTTP masuk lewat
router, ditangani controller, mengambil/mengubah data lewat model, sampai
respons dikembalikan lewat view}
\label{fig:laravel-lifecycle}
\end{figure}

Alur pada \cref{fig:laravel-lifecycle} adalah pola yang akan berulang di
hampir setiap bab buku ini: \route{/pos} pada Bab~3 mengikuti alur ini
persis, begitu juga setiap endpoint API pada Bab~9.

\begin{notebox}
Laravel bukan satu-satunya kerangka kerja PHP, dan PHP bukan satu-satunya
bahasa untuk membangun aplikasi web monolith. Symfony menawarkan fleksibilitas
arsitektur yang lebih tinggi dengan kurva belajar yang lebih curam; Django
(Python) dan Ruby on Rails punya filosofi ``convention over configuration''
yang mirip Laravel. Posisi Laravel yang menonjol ada pada ekosistem
paketnya (Sanctum untuk API, Excel untuk impor/ekspor, Cashier untuk
pembayaran) yang langsung terintegrasi tanpa konfigurasi berlebihan, sesuatu
yang akan terasa langsung mulai Bab~9.
\end{notebox}

\section{Menyiapkan Proyek Laravel dan SQLite}
\label{sec:arsitektur-web-modern-2}

Setiap proyek Laravel butuh basis data untuk berjalan, dan buku ini memilih
SQLite alih-alih MySQL atau PostgreSQL untuk satu alasan praktis: SQLite
menyimpan seluruh basis data dalam satu berkas biasa, tanpa proses server
terpisah yang harus dinyalakan, dikonfigurasi, dan diberi kredensial
sebelum baris kode pertama dijalankan. Pembaca yang baru menginstal PHP
bisa langsung menjalankan migrasi pada menit pertama, bukan menghabiskan
setengah jam mengatur akun MySQL. Simple POS sendiri berjalan di atas
SQLite bahkan untuk demonstrasi kelas, bukan cuma untuk latihan buku ini.

\begin{lstlisting}[language=bash]
composer create-project laravel/laravel simple-pos
cd simple-pos
cp .env.example .env
php artisan key:generate
\end{lstlisting}

Tiga perintah terakhir menyiapkan berkas konfigurasi \file{.env} (tempat
kredensial dan pengaturan lingkungan disimpan, terpisah dari kode sumber
supaya tidak pernah ikut ter-commit ke Git) dan menghasilkan kunci enkripsi
aplikasi yang dipakai Laravel untuk mengamankan sesi dan cookie.

Selain dependensi PHP yang baru saja diinstal \cmd{composer}, proyek Laravel
terbaru juga menyertakan \file{package.json}: daftar dependensi JavaScript
untuk toolchain frontend (Vite dan Tailwind CSS) yang akan mulai dipakai
Bab~3. \cmd{npm install} adalah padanan \cmd{composer install} di dunia
JavaScript, membaca \file{package.json} lalu mengunduh setiap paket yang
terdaftar ke folder \file{node\_modules/}:

\begin{lstlisting}[language=bash]
npm install
\end{lstlisting}

Langkah ini belum wajib untuk menjalankan Simple POS hari ini; halaman
bawaan Laravel masih tampil dengan gaya cadangan meski \file{node\_modules/}
belum ada. Menjalankannya sekarang menghindari jeda instalasi mendadak
begitu Bab~3 mulai memakai Tailwind dan Alpine.js lewat Vite. Langkah
berikutnya mengarahkan koneksi basis data ke sebuah berkas SQLite:

\begin{lstlisting}[language=bash]
touch database/database.sqlite
php artisan migrate:fresh --seed
\end{lstlisting}

\artisan{migrate:fresh} menjalankan seluruh berkas migrasi dari awal
(menghapus tabel lama jika ada, lalu membangunnya kembali dari definisi
skema terbaru), sedangkan opsi \cmd{--seed} langsung mengisi tabel yang baru
dibuat dengan data contoh lewat berkas seeder. Untuk Simple POS, seeder ini
mengisi ratusan produk dan ribuan transaksi sekaligus, sesuatu yang akan
dibahas lebih dalam pada Bab~4 ketika alasan mengisi data dalam skala nyata,
bukan sekadar beberapa baris, mulai terlihat manfaatnya.

\begin{warningbox}
Berkas \file{.env} menyimpan informasi yang sensitif (kunci aplikasi,
kredensial basis data produksi jika ada) dan tidak boleh ikut di-commit ke
Git. Proyek hasil \cmd{composer create-project} sudah menyertakan berkas
\file{.gitignore} yang mengecualikan \file{.env} secara default; jangan
pernah menghapus baris itu.
\end{warningbox}

Setelah migrasi selesai, server pengembangan bawaan Laravel bisa langsung
dijalankan tanpa konfigurasi web server terpisah seperti Apache atau Nginx:

\begin{lstlisting}[language=bash]
php artisan serve
\end{lstlisting}

Perintah ini menjalankan aplikasi pada \route{http://127.0.0.1:8000} memakai
server PHP bawaan, cukup untuk seluruh pengembangan lokal di buku ini.
Server produksi sungguhan (dibahas pada Bab~11) memakai konfigurasi
berbeda, tetapi struktur aplikasi yang dipelajari di sini tetap sama persis.

\section{Praktik: Inisialisasi Simple POS dan Alur Git}
\label{sec:arsitektur-web-modern-3}

\subsection*{Praktik 1.1: Mengenal Struktur Proyek}

Sebelum menulis kode fitur pertama, penting mengenali bagian mana dari
struktur folder Laravel yang akan sering disentuh sepanjang buku ini, dan
mana yang boleh diabaikan untuk sementara.

\begin{enumerate}
  \item Buka folder \file{routes/}, khususnya \file{routes/web.php}: di
    sinilah setiap alamat URL yang bisa diakses pengguna aplikasi
    didaftarkan.
  \item Buka folder \file{app/Http/Controllers/}: tempat class yang
    memproses setiap permintaan (menerima input, memanggil model, memilih
    tampilan) akan bertambah satu demi satu mulai Bab~2.
  \item Buka folder \file{database/migrations/}: tempat perubahan skema
    basis data didefinisikan sebagai kode, bukan diklik lewat aplikasi
    basis data terpisah.
  \item Bandingkan dengan folder \file{vendor/}: berisi seluruh paket
    Composer pihak ketiga, tidak pernah diedit manual, dan tidak ikut
    di-commit ke Git.
\end{enumerate}

Struktur ini bukan kebetulan. Ia mewujudkan pola MVC yang disinggung di
istilah penting bagian awal bab: \file{routes/} dan
\file{app/Http/Controllers/} menangani sisi Controller, model Eloquent yang
akan diperkenalkan di Bab~5 menangani sisi Model, dan berkas Blade di
\file{resources/views/} (dibahas mulai Bab~3) menangani sisi View.

\subsection*{Praktik 1.2: Konvensi Commit Increment pada Simple POS}

Repositori Simple POS yang menjadi rujukan buku ini tidak dibangun sekaligus
sebagai satu proyek utuh, melainkan tumbuh commit demi commit, masing-masing
diberi label \cmd{increment N}. Increment 1 hanya berisi kerangka proyek
kosong hasil \cmd{composer create-project}, persis seperti hasil Praktik
1.1; increment berikutnya menambah satu kemampuan konkret, routing dan
controller produk, lalu tampilan Blade, lalu skema basis data, mengikuti
urutan yang sama dengan urutan bab pada buku ini. Membaca riwayat commit
Simple POS secara berurutan sama seperti membaca buku ini bab demi bab,
hanya dalam bentuk kode sungguhan alih-alih prosa.

\begin{tipbox}
Kode lengkap setiap bab tersedia dalam dua bentuk. Repositori
\texttt{\repobase} menyimpan seluruh riwayatnya dalam satu tempat, dengan
satu branch per bab (\texttt{chapter-01} hingga \texttt{chapter-12}), setiap
branch dibangun di atas branch bab sebelumnya; cocok dibaca berurutan
seperti membaca riwayat commit sungguhan. Setiap akhir bagian Praktik pada
buku ini, seperti kotak di bawah, sebaliknya menunjuk ke repositori
template tersendiri untuk bab itu saja (mis. \texttt{\repoowner/simple-pos-ch04}
untuk Bab~4): klik tombol \textbf{Use this template} di GitHub untuk
langsung mulai dari kondisi bab itu, tanpa perlu tahu cara \cmd{git
checkout} branch tertentu. Buku ini menampilkan potongan kode yang paling
penting untuk dipahami; kode penuh yang siap dijalankan selalu ada di
kedua repositori itu.
\end{tipbox}

\branchref{chapter-01}

\section{Rangkuman}
\label{sec:arsitektur-web-modern-rangkuman}

\begin{itemize}
  \item Monolith menyatukan seluruh lapisan aplikasi dalam satu basis kode
    dan satu proses deploy, cocok untuk skala Simple POS; microservices
    memecahnya menjadi layanan independen dengan ongkos operasional yang
    hanya sepadan untuk sistem berskala besar; serverless cocok untuk beban
    kerja sporadis, bukan aplikasi dengan koneksi basis data yang konsisten.
  \item Proyek Laravel 13 disiapkan dengan \cmd{composer create-project},
    dihubungkan ke SQLite sebagai basis data zero-setup, lalu dijalankan
    dengan \artisan{migrate:fresh --seed} dan \artisan{serve}.
  \item Konvensi commit \cmd{increment N} pada repositori Simple POS
    memetakan pertumbuhan aplikasi bab demi bab, dengan kode lengkap tiap
    bab tersedia sebagai branch tersendiri di GitHub.
\end{itemize}

\section{Latihan}

\begin{enumerate}
  \item Jalankan ketiga perintah pada Praktik 1.1 di komputermu sendiri,
    lalu verifikasi bahwa \cmd{php artisan serve} menampilkan halaman
    selamat datang Laravel di \route{http://127.0.0.1:8000}.
  \item Buka \file{database/migrations/} pada proyek yang baru dibuat, lalu
    identifikasi berapa banyak berkas migrasi yang sudah ada secara bawaan
    (tanpa menulis kode apa pun) dan tabel apa saja yang dibuatnya.
  \item Jelaskan dengan kata-katamu sendiri: mengapa sebuah warung dengan
    satu kasir lebih cocok memakai arsitektur monolith dibanding
    microservices, sementara sebuah platform e-commerce nasional dengan
    jutaan pengguna sering memakai microservices.
  \item \textbf{Evaluasi mandiri:} tanpa membuka kembali isi bab ini, coba
    jelaskan dengan kata-katamu sendiri: (a) perbedaan monolith dan
    microservices, (b) mengapa SQLite dipilih untuk buku ini, dan (c) apa
    fungsi opsi \cmd{--seed} pada \artisan{migrate:fresh}. Periksa kembali
    jawabanmu dengan isi bab ini, dan catat bagian mana yang masih perlu
    dipelajari ulang.
\end{enumerate}

\section{Referensi}

Pembahasan pada bab ini merujuk pada dokumentasi resmi Laravel
\cite{laraveldocs} untuk struktur proyek dan perintah Artisan, serta
Manual PHP \cite{phpmanual} untuk rujukan bahasa yang mendasari kerangka
kerja Laravel.
